
\section{conclusion}

\begin{frame}
\frametitle{Strong-Field Ionization Phenomena Revealed by Quantum Trajectories \\
\small (Phys. Rev. Lett. \textbf{133}, 073201 (2024))}

\framesubtitle{Closure}

\small
\begin{textblock*}{400pt}(30pt,60pt)
\begin{itemize}
\footnotesize
\item Quantum trajectories represent a powerful alternative tool
to investigate electron dynamics beneath and above the potential barrier 
\item The tunneling parameter $\bar{\beta}$ allows a mapping 
of the angular-resolved tunneling dynamics 
\item The quantum adiabatic parameter$\gamma_q$ is superior to 
the Keldysh parameter $\gamma_{\mbox{\tiny$\rm K$}}$ in the sense that
\begin{itemize}
\item[$\bullet$] Rare quantitative information can be probed about 
various high intensity effects
\item[$\bullet$] Importance of long range electronic potential 
and nonadiabatic effects are better rendered 
\end{itemize}

\vspace{20pt}

\item This might pave the way to a  better understanding of 
other routinely studied dynamical effects (HHG, etc..)
\end{itemize}

\end{textblock*}


%\begin{textblock*}{100pt}(260pt,70pt)
%\includegraphics[scale=.30]{images/medium1.png}
%\end{textblock*}

\end{frame}
